% !PRINTED 44
\weberpage{44}

%<<CONT
representable, so that
\weq{3}{0 \leqq \alpha_1 < a_1,\ 0 \leqq \alpha_2 < a_2, \ldots 0 \leqq \alpha_\nu < a_\nu.}

The exponents $\alpha_1, \alpha_2 \ldots \alpha_\nu$, or any numbers congruent to them
with respect to the moduli $a_1, a_2 \ldots a_\nu$, are called the \gesperrt{indices}
of the element $A$, and the following theorem holds:

1. \gesperrt{One obtains the indices of a compositum $A A'$
by adding the corresponding indices of the two
factors.}

If we assign to each element $A$ in any manner a \gesperrt{numerical value},
then we may call this assignment a function
of $A$. Such a function $\chi(A)$ shall now be called a
\gesperrt{group character} if for every two elements
$A$, $A'$ of $S$ it satisfies the condition:
\weq{4}{\chi(AA') = \chi(A)\chi(A'),}
and consequently also the more general one
\[ \chi(AA'A''\ldots) = \chi(A)\chi(A')\chi(A'')\ldots \]
Two such functions $\chi$ and $\chi_1$ are regarded as different
if there is at least one element $A$ for which $\chi(A)$
is different from $\chi_1(A)$. We shall now determine these
characters more closely.

If we set $A' = 1$, then from (4) there results:
\weq{5}{\chi(A) = \chi(A)\chi(1),}
and if we exclude once and for all the case that all
$\chi(A) = 0$:
\weq{6}{\chi(1) = 1,}
hence the first theorem:

2. \gesperrt{Every character has the value 1 for the unit
element.}

If we further set $A' = A$, then from (4) it follows:
\[ \chi(A^2) = [\chi(A)]^2, \]
and from this by complete induction for every exponent $h$:
\weq{7}{\chi(A^h) = \chi(A)^h.}
If we therefore denote by $a$ the degree of the element $A$, then
it follows, when one sets $h = a$ in (7) and uses (6):
\weq{8}{\chi(A)^a = 1,}
hence:
%CONT>>
